\chapter{Mode-I Benchmark and Fixed-Mesh Reference}
\label{ch:mode1_reference}

\section{Benchmark Geometry and Sharp Seam Representation}
\label{sec:benchmark_geometry}

The foundational benchmark is a $1.0\,\mathrm{mm} \times 1.0\,\mathrm{mm}$ square plate containing an initial edge crack of length $a_0 = 0.5\,\mathrm{mm}$ located along the horizontal symmetry line $y = 0.5\,\mathrm{mm}$ (Fig.~\ref{fig:mode1_geometry}). The crack is modeled as a zero-gap sharp seam with duplicate, coincident nodal pairs rather than a finite-width notch.

\begin{figure}[htbp]
  \centering
  \includegraphics[width=0.52\textwidth]{fig_mode1_geometry.png}
  \caption{Mode-I benchmark boundary value problem: $1.0\,\mathrm{mm} \times 1.0\,\mathrm{mm}$ square domain with a sharp edge crack of length $a_0 = 0.5\,\mathrm{mm}$ at $y = 0.5\,\mathrm{mm}$. Boundary conditions prescribe roller support ($u_y = 0$) along the bottom edge with a pinned point ($u_x = 0$) to eliminate rigid-body translation, and uniform tensile displacement $u_y = u$ on the top boundary.}
  \label{fig:mode1_geometry}
\end{figure}

\begin{table}[htbp]
\centering
\caption{Material and phase-field constitutive parameters for the Mode-I benchmark \citep{pandey2025}.}
\label{tab:mode1_parameters}
\begin{tabular}{llrl}
\toprule
Parameter & Symbol & Value & Units / Description \\
\midrule
Young's modulus & $E$ & $210.0$ & $\mathrm{GPa}$ ($\mathrm{kN/mm^2}$) \\
Poisson's ratio & $\nu$ & $0.30$ & -- \\
Critical fracture energy release rate & $G_c$ & $2.7 \times 10^{-3}$ & $\mathrm{kN/mm}$ ($2.7\,\mathrm{N/mm}$) \\
Phase-field regularization length scale & $l_0$ & $0.0075$ & $\mathrm{mm}$ ($7.5\,\mu\mathrm{m}$) \\
Residual stiffness parameter & $k_{\mathrm{res}}$ & $1.0 \times 10^{-7}$ & Artificial residual stiffness \\
\bottomrule
\end{tabular}
\end{table}

\subsection{Boundary-Condition Verification and Seam Compliance}
\label{sec:notch_vs_seam}

A vital physical distinction must be drawn between the material Young's modulus $E = 210.0\,\mathrm{GPa}$ and the global initial structural stiffness $K_0$:
\begin{equation}
  K_0 = \left. \frac{\mathrm{d}F}{\mathrm{d}u} \right|_{\text{initial elastic range}} \approx \frac{\Delta F}{\Delta u}.
  \label{eq:k0_definition}
\end{equation}
During initial model verification, representing the crack as a finite-width notch was found to introduce substantial artificial compliance ($K_0 \approx 75.47\,\mathrm{kN/mm}$). The geometry was therefore corrected to the required zero-gap sharp seam before the governed convergence and adaptive-remeshing studies were performed. All subsequent results use the corrected sharp-seam geometry.

\section{Boundary-Condition Specification and Reconstruction Provenance}
\label{sec:boundary_conditions}

The boundary conditions governing the Mode-I benchmark are:
\begin{itemize}
  \item \textbf{Bottom Boundary ($y = 0.0\,\mathrm{mm}$):} Vertical displacement is fixed ($u_y = 0.0$) across all bottom nodes ($\mathtt{N\_BOTTOM}$). A single reference point or midpoint node is restrained in the horizontal direction ($u_x = 0.0$) to eliminate rigid-body motion without over-constraining Poisson contraction.
  \item \textbf{Top Boundary ($y = 1.0\,\mathrm{mm}$):} Prescribed uniform tensile displacement loading $u_y = u(t)$ applied monotonically.
  \item \textbf{Crack Flanks and Lateral Edges ($x = 0, 1.0\,\mathrm{mm}$):} Traction-free boundaries with natural boundary conditions for the phase field $\nabla d \cdot \mathbf{n} = 0$.
\end{itemize}

\section{Verified Fixed-Mesh Reference Anchor}
\label{sec:fixed_reference_anchor}

The fixed-mesh reference solution, Job \texttt{1406015.mmaster02} ($S_1$, $15{,}192$ finite elements, $15{,}522$ total node entries), serves as the authoritative quantitative anchor against which all adaptive remeshing results are evaluated. Under quasi-static displacement-controlled tension, the specimen exhibits linear elasticity up to damage localization, reaching a peak reaction force $F_{\max} = 0.757778\,\mathrm{kN}$ at $u(F_{\max}) = 0.005857\,\mathrm{mm}$, followed by sharp post-peak softening as the crack propagates horizontally through the remaining ligament.

Initial structural stiffness is extracted using an unconstrained ordinary-least-squares (OLS) linear fit over the initial $N = 400$ active increments ($u \le 1.0\,\mu\mathrm{m}$):
\begin{equation}
  F_i = K_0 u_i + b + \epsilon_i, \qquad \boxed{K_0 = 137.945520\,\mathrm{kN/mm}},
  \label{eq:k0_ols_anchor}
\end{equation}
with intercept $b = 4.472368 \times 10^{-5}\,\mathrm{kN}$ and coefficient of determination $R^2 = 0.99999960$. Digitized values from \citet{pandey2025} (Fig.~7(a)) report $F_{\max} \approx 0.758\,\mathrm{kN}$ and $u(F_{\max}) \approx 0.005860\,\mathrm{mm}$, demonstrating excellent numerical agreement (Table~\ref{tab:fixed_reference}). Historical Job \texttt{1398090.mmaster02} is preserved as documented baseline evidence.

\begin{table}[htbp]
\centering
\caption{Verified quantitative metrics for the fixed-mesh reference anchor (Job \texttt{1406015}).}
\label{tab:fixed_reference}
\begin{tabular}{llrl}
\toprule
Quantity & Symbol & Value & Source / Extraction Method \\
\midrule
Finite element count & $N_{\mathrm{mesh}}$ & $15{,}192$ & Structured CPE4 mesh \\
Total node entries & $N_{\mathrm{nodes}}$ & $15{,}522$ & Layer 1/2/3 coincident nodes + RP \\
Peak reaction force & $F_{\max}$ & $0.757778\,\mathrm{kN}$ & Sum of top reaction forces $R_{y}$ \\
Peak displacement & $u(F_{\max})$ & $0.005857\,\mathrm{mm}$ & Prescribed stroke at limit load \\
Initial structural stiffness & $K_0$ & $137.945520\,\mathrm{kN/mm}$ & OLS fit ($N=400$, $u \le 1.0\,\mu\mathrm{m}$) \\
OLS intercept & $b$ & $4.472368 \times 10^{-5}\,\mathrm{kN}$ & Intercept from linear regression \\
Coefficient of determination & $R^2$ & $0.99999960$ & Goodness of fit \\
\bottomrule
\end{tabular}
\end{table}

\section{Resolution of the Boundary-Set Formatting Defect}
\label{sec:nbottom_defect}

During initial deployment of the native adaptive 71,320-element mesh (historical Job \texttt{1399632}), an apparent stiffness deficit was observed ($K_0 \approx 122.38\,\mathrm{kN/mm}$, $-11.28\%$). Forensic audit revealed that this was caused by an Abaqus keyword input card limit: the node set \texttt{N\_BOTTOM} contained 150 node identifiers formatted on a single data line without \texttt{GENERATE}. Abaqus \texttt{pre} silently parsed only the first 16 node identifiers and omitted the remaining 134 entries.

Wrapping the node identifiers across valid data lines ($\le 16$ entries per card) restored all 150 vertical boundary constraints. The defect was independently verified and is classified as \textbf{\texttt{RESOLVED\_AND\_CLOSED}}; all subsequent adaptive results use the corrected boundary-condition definition.

The corrected input deck was verified in authoritative full-fracture Job \texttt{1404933} and frozen elastic Job \texttt{1405044}. As shown in Table~\ref{tab:nbottom_comparison} and Fig.~\ref{fig:mode1_fu_recovery}, proper boundary set wrapping restores initial stiffness to $K_0 = 137.820804\,\mathrm{kN/mm}$ ($-0.09\%$ vs reference anchor), definitively resolving this diagnostic issue.

\begin{table}[htbp]
\centering
\caption{Mechanical response recovery after resolving the \texttt{N\_BOTTOM} keyword formatting defect.}
\label{tab:nbottom_comparison}
\small
\begin{tabularx}{\textwidth}{lrrrrY}
\toprule
Case Description & Job ID & Elements & $K_0$ [kN/mm] & $F_{\max}$ [kN] & Status / Governed Role \\
\midrule
Fixed Reference Anchor ($S_1$) & \texttt{1406015} & $15{,}192$ & $137.945520$ & $0.757778$ & Canonical reference anchor \\
Corrected Full Fracture ($A_1$) & \texttt{1404933} & $71{,}320$ & $137.820804$ & $0.745325$ & Production $1\%$ adaptive mesh \\
Frozen Requalification & \texttt{1405044} & $71{,}320$ & $138.021013$ & -- & Independent elastic verification \\
\bottomrule
\end{tabularx}
\end{table}

\begin{figure}[htbp]
  \centering
  \includegraphics[width=0.86\textwidth]{fig_mode1_fu_recovery.pdf}
  \caption{Force--displacement comparison between the canonical fixed-mesh $S_1$ reference ($15{,}192$ finite elements) and the corrected native-adaptive $A_1$ model with \texttt{errorTarget = 1\%} ($71{,}320$ finite elements). The initial structural stiffness differs by only $0.09\%$, while the adaptive peak reaction force is approximately $1.64\%$ below the reference value.}
  \label{fig:mode1_fu_recovery}
\end{figure}
